\par\clearpage\SourceLeaf{310}\NoSourceHead
\HistoricalFoot{X\textsubscript{2} Abh. 3.}{1}
\begin{center}
{\Large\bfseries ÜBER DEN ERSTEN UND VIERTEN\\
GAUSSSCHEN BEWEIS DES FUNDAMENTAL-\\
SATZES DER ALGEBRA\par}
\vspace{2cm}
VON
\vspace{2cm}

{\large\bfseries ALEXANDER OSTROWSKI\par}
\vfill
\rule{\linewidth}{.4pt}\par
{\small Neubearbeitung des 1920 als Anhang zum Heft VIII der \textit{Materialien für eine wissenschaftliche Biographie von Gauss} (S. 50—58) erschienenen Aufsatzes.\par
Nachrichten der K. Gesellschaft der Wissenschaften zu Göttingen. Mathem.-physik. Klasse. 1920 Beiheft.\par}
\rule{\linewidth}{.4pt}
\end{center}
\clearpage\SourceLeaf{312}\NoSourceHead\HistoricalFoot{}{1\textsuperscript{*}}
\renewcommand{\thefootnote}{\arabic{footnote})}\setcounter{footnote}{0}
\begin{center}\textbf{Einleitung.}\end{center}
\noindent\hspace*{\parindent}Während die im ersten Teil der \textsc{Gauss}schen \textit{Dissertation} (1799, Werke
III, S. 1) enthaltene Besprechung der früheren Beweisversuche des Fundamentalsatzes
der Algebra sich durch ganz ausserordentliche Sorgfalt auszeichnet, fällt
daneben der im zweiten Teil entwickelte Beweis dieses Satzes etwas ab. Nicht
etwa, weil dieser Beweis in geometrischer Einkleidung vorgetragen wird, sondern,
weil bei ihm Eigenschaften der algebraischen Kurven verwendet werden,
die weder in der \textit{Dissertation} selbst, noch in der vorgaussschen Literatur bewiesen
sind.
\par
In den folgenden Ausführungen soll nun versucht werden, für den ersten
und vierten \textsc{Gauss}schen Beweis des Fundamentalsatzes der Algebra eine vollständig
strenge Darstellung zu geben.
\par
Selbstverständlich darf ein Beweis eines analytischen Satzes, der geometrische
Begriffe benutzt, deshalb allein noch nicht als unstreng bezeichnet
werden. Solange die geometrische Anschauung nur dazu dient, den Überblick
über die benutzten Voraussetzungen und den Gedankengang des Beweises zu
erleichtern — oder auch dazu, direkte Anwendung und Formulierung der Grundeigenschaften
der Zahlenreihe und der stetigen Funktionen aus darstellungstechnischen
Gründen zu umgehen — werden wir heute einen solchen Beweis
als korrekt bezeichnen, da man sich eben auf den Standpunkt stellen kann,
dass die inzwischen erfolgte sorgfältige Grundlegung der Analysis das an sich
etwas schwankende Fundament aller solchen Beweise mit einem Schlage gesichert
hat.
\par
Man könnte sich natürlich auch auf den intransigenten Standpunkt stellen,
dass, solange keine Klarheit über den Begriff der Irrationalzahl bestand, ein
strenger Beweis eines derartigen Existenzsatzes überhaupt unmöglich war. Doch
dürfte es wohl richtiger sein, einen Beweis der »vorkritischen« Periode unserer\par
\clearpage\SourceLeaf{313}\ArithmeticHead{4}{L}{A. OSTROWSKI, ÜBER DEN ERSTEN UND VIERTEN GAUSSSCHEN BEWEIS USW.}
\noindent Wissenschaft als streng anzusehen, wenn er sich an Hand der inzwischen aufgestellten
Definitionen und einfachsten Elementarsätze ohne weiteres ergänzen
lässt. In diesem Sinne darf der zweite \textsc{Gauss}sche Beweis (Werke III, S. 31)
als absolut streng bezeichnet werden.
\par
Anders ist es, wenn die geometrische Einkleidung dazu dient, die wirklichen
Schwierigkeiten zu verhüllen und in die Daten des Problems neue,
nicht ausdrücklich formulierte Voraussetzungen hineinzutragen. In diesem
Falle darf man durchaus von einer wesentlichen Lücke sprechen. —
\par
In dem \textsc{Gauss}schen Beweise handelt es sich im wesentlichen um folgendes:
Es wird ein Polynom
\[
f(z)=z^n+Az^{n-1}+\cdots+C
\]
mit reellen Koeffizienten betrachtet, und es werden die beiden Kurven
\[
\text{(I)}\quad\Re f(z)=0,\qquad\text{(II)}\quad\Im f(z)=0
\]
untersucht. Es gibt dann ein \(R\) derart, dass auf jedem der Kreise \(|z|=r\),
\(r\geqq R\) genau \(2n\) Punkte der Kurve (I) und ebensoviele Punkte der Kurve (II)
liegen, wobei die beiden Punktsysteme einander trennen. Benutzt man Polarkoordinaten
\(r,\varphi\), so lassen sich die Punkte von (I) und (II) für \(r\geqq R\) zu je
\(2n\) Kurvenzügen
\[
\varphi=A_\mu(r),\quad\varphi=B_\mu(r),\qquad(\mu=1,2,\ldots,2n)
\]
zusammenfügen, wobei, wie \textsc{Gauss} im Art. 20 bemerkt, sich leicht beweisen
lässt, dass die Funktionen \(A_\mu,B_\mu\) in \(r\) stetig sind. Die zugehörigen Kurvenzüge
erstrecken sich von der Kreislinie \(|z|=R\) ausgehend ins Unendliche,
und zwar so, dass zwischen je zwei aufeinanderfolgenden Zweigen der Kurve (I)
je ein Zweig der Kurve (II) liegt und umgekehrt. Überdies ergibt sich
aus dem Beweise der obigen Tatsachen, dass \(\Re f(z)\) beim Überschreiten der
Kurve (I) jedesmal sein Vorzeichen wechselt, während in den zwischen je zwei
aufeinanderfolgenden Zweigen der Kurve (I) gelegenen Bereichen und auf den
entsprechenden Bogen von \(|z|=R\) das Vorzeichen von \(\Re f(z)\) jedesmal unverändert
bleibt. Das Analoge gilt für die Kurve (II) und \(\Im f(z)\). — Für die
Beweise dieser Tatsachen vergl. die Nummern 1 und 2 der folgenden Darstellung.
— Es handelt sich nunmehr um den Verlauf der Kurven (I), (II) im
Innern von \(|z|=R\), und die \textsc{Gauss}sche Betrachtung beruht in erster Linie\par
\clearpage\SourceLeaf{314}\ArithmeticHead{5}{R}{A. OSTROWSKI, ÜBER DEN ERSTEN UND VIERTEN GAUSSSCHEN BEWEIS USW.}
\noindent auf der Tatsache, dass sich die ausserhalb von \(|z|=R\) verlaufenden Zweige
von (I) (und ebenso auch diejenigen von (II)) so zu Paaren zusammenfassen
lassen, dass, wenn man längs des einen Zweiges \(\varphi=A_\mu(r)\) eines solchen
Paares ins Innere von \(|z|=R\) hineingeht, man die ganze Kreisfläche durchquert
und schliesslich längs des anderen Zweiges \(\varphi=A_\nu(r)\) dieses Paares aus
dem Kreise heraustritt. Im Nachweis dieser Tatsache liegt die Hauptschwierigkeit;
die weiteren Schlüsse der \textsc{Gauss}schen \textit{Dissertation} folgen dann in durchaus
einwandfreier Weise.
\par
Zur Begründung der genannten Tatsache sagt nun \textsc{Gauss}, es sei aus der
höheren Geometrie bekannt, dass, wenn ein Zweig einer algebraischen Kurve
in einen begrenzten Raum eintritt, er notwendig wieder aus demselben heraustreten
müsse, und er bemerkt zur Erläuterung in einer Fussnote (Werke III,
S. 27), dass eine algebraische Kurve »weder plötzlich abbricht, noch sich nach
unendlich vielen Umläufen gewissermassen in einem Punkt verlieren kann«.
Er fügt dann hinzu: »\mbox{.\,.\,.} und soviel ich weiss, hat noch Niemand hiergegen
einen Zweifel vorgebracht. Doch werde ich, wenn es jemand fordert, bei
anderer Gelegenheit unternehmen, einen keinem Zweifel unterworfenen Beweis
zu liefern.« Weiter sagt \textsc{Gauss}, dass im vorliegenden Falle das Nichtabbrechen
der hier vorkommenden Kurven auch schon daraus hervorgeht, dass
sie Gebiete verschiedenen Vorzeichens von \(\Re f(z)\) bezw. \(\Im f(z)\) voneinander
trennen.
\par
Hierzu ist nun zu bemerken, dass die genaue Durchführung der zuletzt
angegebenen Überlegung einer sorgfältigen topologischen Ausführung bedarf.
Vor allem wäre aber auch der Beweis, dass ein solches Verhalten unmöglich
ist, insofern nicht genügend, als daraus die Möglichkeit des »Fortschreitens«
längs der im Innern von \(|z|=R\) liegenden Zweige unserer Kurven, d. h. die
Existenz einer Parameterdarstellung, noch nicht folgt, wie denn überhaupt
der Beweis für die Existenz einer Parameterdarstellung \textit{durch die Fortsetzung
im Kleinen} unter Umständen auf grosse Schwierigkeiten stossen kann, falls
man als einen solchen Parameter nicht eine der Koordinaten oder die Bogenlänge
der Kurve benutzen kann. Und die Tatsache, dass es sich bei den
Zweigen einer algebraischen Kurve um einfache Linien handelt, die sich also
etwa aus endlich vielen regulären Bogen zusammensetzen, darf selbstverständlich\par
\clearpage\SourceLeaf{315}\ArithmeticHead{6}{L}{A. OSTROWSKI, ÜBER DEN ERSTEN UND VIERTEN GAUSSSCHEN BEWEIS USW.}
\noindent nicht den Erfahrungen entnommen werden, die man bei der Untersuchung
von Kurven etwa der ersten sechs Ordnungen gesammelt hat.
\par
Demnach hat man zur Ergänzung des ersten \textsc{Gauss}schen Beweises zu zeigen,
dass man jede algebraische Kurve — oder auch nur die im \textsc{Gauss}schen Beweis
vorkommenden Kurven — in endlich viele reguläre Bogen zerlegen kann,
die sich durch eine Gleichung von einer der beiden Formen
\[
y=g(x),\quad x=g(y)
\]
darstellen lassen, wobei die Funktion \(g\) im Innern des betrachteten Intervalls
stetig und beschränkt differenzierbar ist.
\par
Wir führen nun in der weiter unten folgenden Darstellung den Beweis
für die \textsc{Gauss}schen Kurven (I), (II), obgleich er, wie wir bemerken möchten,
auch für beliebige algebraische und sogar für nur analytische Kurven ausreicht.
Die Überlegungen, die dabei zu benutzen sind, sind zunächst algebraischer
Natur. Es ist nämlich für die Durchführung unseres Beweises sehr
wertvoll, dass man von vornherein die endlich vielen Punkte angeben kann,
in denen ein regulärer Teilbogen der betrachteten Kurve enden könnte.
Hierzu braucht man offenbar in erster Linie nur die singulären Punkte und
die Berührungspunkte der zur \(x\)- und \(y\)-Achse parallelen Tangenten der betrachteten
Kurve aufzusuchen. Es ist daher wichtig, die Gleichung der Kurve
so einzurichten, dass solche Punkte nur in endlicher Anzahl vorkommen.
Trifft dies aber in einem Falle nicht zu, so folgt daraus, geometrisch gesprochen,
dass unsere Kurve ein mehrfaches Stück enthält, und es ist wesentlich,
derartige, mehrfach auftretende Stücke durch algebraische Operationen
entfernen zu können. Dies wird möglich, wenn man die folgende algebraische
Tatsache benutzt:
\par
\textls[100]{Hat ein Polynom in zwei Veränderlichen unendlich viele gemeinsame
Nullstellen mit seiner Ableitung nach einer von diesen
Veränderlichen, so\linebreak besitzt es einen mehrfachen Faktor.}
\par
Diese Tatsache ist leicht zu beweisen, z. B. mit Hilfe des \textsc{Euklidischen}
Algorithmus, wobei dann allerdings die folgende für alle Sätze über die Teilbarkeit
von Polynomen fundamentale Erweiterung des \textsc{Gauss}schen Satzes über
die Teiler ganzzahliger Polynome (\textit{Disqu. arithm.} Art. 42, Werke I, S. 34)
heranzuziehen ist:
\par\clearpage\SourceLeaf{316}\ArithmeticHead{7}{R}{A. OSTROWSKI, ÜBER DEN ERSTEN UND VIERTEN GAUSSSCHEN BEWEIS USW.}
\textls[100]{Multipliziert man zwei Polynome}
\[
\begin{aligned}
F(x,y)&=A_0(x)y^m+A_1(x)y^{m-1}+\cdots+A_m(x),\\
G(x,y)&=B_0(x)y^n+B_1(x)y^{n-1}+\cdots+B_n(x),
\end{aligned}
\]
\textls[100]{wo \(A_0,\ldots,A_m;\ B_0,\ldots,B_n\) Polynome in \(x\) sind und sowohl der grösste\linebreak
gemeinsame Teiler von \(A_0,\ldots,A_m\) als auch der grösste gemeinsame\linebreak
Teiler von \(B_0,\ldots,B_n\) gleich 1 ist, so ist im Produkt}
\[
FG=C_0(x)y^{m+n}+C_1(x)y^{m+n-1}+\cdots+C_{m+n}(x)
\]
\textls[100]{der grösste gemeinsame Teiler der Polynome \(C_0,\ldots,C_{m+n}\) auch
gleich 1.}
\par
Es ist von Interesse zu bemerken, dass \textsc{Gauss} den erwähnten arithmetisch-algebraischen
Satz am 22. Juli 1797 gefunden hat, also einige Monate vor der
Entdeckung der Prinzipien des ersten Beweises des Fundamentalsatzes. Daher
ist es durchaus gerechtfertigt, anzunehmen, dass die zur Durchführung einer
Diskussion des Verlaufes einer algebraischen Kurve im Grossen erforderlichen
Hilfsmittel \textsc{Gauss} bei der Abfassung seiner \textit{Dissertation} zur Verfügung standen.
\par
Was den Verlauf eines einzelnen Teilbogens der Kurve zwischen den
auf der Kurve markierten Punkten anbetrifft, so kann man hier die Fortsetzbarkeit
eines solchen Bogens ohne weiteres durch eine Überlegung beweisen,
wie sie offenbar \textsc{Gauss} im Art. 20 der \textit{Dissertation} zum Beweise der Stetigkeit
der dort betrachteten Kurvenzweige vorgeschwebt hat. Um allerdings solche
»infinitesimale« Schritte zu einem vollständigen Teilbogen zusammenzufügen,
muss man wohl von der gleichmässigen Stetigkeit eines Polynoms in zwei
Variabeln Gebrauch machen. Die gleichmässige Stetigkeit lässt sich aber im
Falle von Polynomen sehr einfach beweisen, da man die Differenz zweier
Polynomwerte ohne weiteres abschätzen kann.
\par
Insofern sind die Hilfsmittel, die wir im Folgenden benutzen, alle der
Art, dass sie auch \textsc{Gauss} im Prinzip bekannt gewesen sein mögen. Gerade
in der Algebra und in der Analysis situs lässt sich aber sehr oft ein an sich
einfacher Gedanke nur mit grosser Umständlichkeit durchführen, und wenn
man bedenkt, dass \textsc{Gauss} sich bei einer solchen Durchführung nirgends an
irgendwelche Vorbilder hätte halten können, so wird man verstehen, warum
er sich für eine, für unsere Begriffe reichlich summarische Darstellung entschlossen
hat. Denn schliesslich muss man auch bei einer solchen Untersuchung\par
\clearpage\SourceLeaf{317}\ArithmeticHead{8}{L}{A. OSTROWSKI, ÜBER DEN ERSTEN UND VIERTEN GAUSSSCHEN BEWEIS USW.}
\noindent irgendwo anfangen, und eine vollständige Neubegründung der Elemente
aus einer derartigen Veranlassung zu geben, dürfte wohl \textsc{Gauss} fern
gelegen haben. Andererseits ist er auf die oben gekennzeichnete Lücke auch
später nie zurückgekommen.
\par
In der fünfzig Jahre später entstandenen Neufassung des Beweises von
1799 (in der sogenannten \textit{Jubiläumsschrift} von 1849, Werke III, S. 71) hat
er vielmehr den Ansatz etwas zu modifizieren versucht, anscheinend, um jener
Schwierigkeit aus dem Wege zu gehen. Er betrachtet jetzt nicht mehr die
einzelnen Zweige der Kurven \(\Re f(z)=0,\ \Im f(z)=0\), sondern die zwischen
diesen Zweigen liegenden zusammenhängenden Gebiete, in denen die Funktionen
\(\Re f(z),\ \Im f(z)\) konstante Vorzeichen haben. Allerdings gelingt es ihm
auch so nicht, die Betrachtung der im Innern von \(|z|=R\) liegenden Zweige
der Kurven (I), (II) zu vermeiden, er betrachtet vielmehr auch in der Neufassung
seines Beweises Umläufe längs dieser Linien, wozu wiederum eine
vorherige Diskussion des geometrischen Charakters dieser Linien erforderlich
wäre. Der einzige Gewinn scheint darin zu bestehen, dass man sich nunmehr
auf die Betrachtung nur einer der beiden Kurven (I), (II) beschränken
kann und einen besseren Überblick über die \textls[100]{Gesamtheit} der Nullstellen
von \(f(z)\) hat.
\par
Andererseits erweist sich der Gedanke, anstatt der Linien \(\Re f(z)=0\),
\(\Im f(z)=0\) die Gebiete \(\Re f(z)\gtrless0,\ \Im f(z)\gtrless0\) zu betrachten, von wesentlich
grösserer Tragweite, als man danach hätte annehmen können. Wir erinnern
daran, dass in der Theorie der konformen Abbildung und in
den neueren Untersuchungen der mengentheoretischen Topologie der Gedanke,
anstatt einer Jordankurve die von ihr begrenzten Gebiete zu betrachten,
sich als sehr fruchtbar erwiesen und z. B. zu einer wesentlichen Vereinfachung
des Beweises des \textsc{Jordan}schen Kurvensatzes geführt hat. Es ist
daher sicher von Interesse, dass man, von der Betrachtung der \textsc{Gauss}schen
Flächenstücke ausgehend, den Beweis des Fundamentalsatzes so zu Ende
führen kann, dass dabei ein weiteres Eingehen auf die Gestalt der im Innern
von \(|z|\leqq R\) liegenden Teile der Kurven (I), (II) vollständig überflüssig wird.
Man hat dabei nur mit dem Begriff eines Kontinuums zu arbeiten und die verschiedenen
Randkontinua der \textsc{Gauss}schen Flächenstücke zu untersuchen. Dies
erscheint umso bemerkenswerter, als dabei von den speziellen Eigenschaften\par
\clearpage\SourceLeaf{318}\ArithmeticHead{9}{R}{A. OSTROWSKI, ÜBER DEN ERSTEN UND VIERTEN GAUSSSCHEN BEWEIS USW.}
\HistoricalFoot{X\textsubscript{2} Abh. 3.}{2}
\noindent der stetigen Funktionen \(\Re f(z),\ \Im f(z)\) im Innern des Kreises \(|z|=R\) kein Gebrauch
gemacht und so der folgende äusserst allgemeine Satz gewonnen wird:
\par
\textls[100]{Haben zwei im Innern und auf dem Rande von \(|z|=R\) stetige\linebreak
reelle Funktionen \(F,G\) auf der Kreislinie \(|z|=R\) die gleiche Anzahl\linebreak
von einander trennenden Vorzeichenwechseln, und verschwinden\linebreak
sie sonst auf der Kreislinie nicht, so haben sie im
Innern des Kreises wenigstens eine gemeinsame Nullstelle.} — Hierin
liegt eine wohl unerwartet direkte Bestätigung der Bemerkungen am Schlusse
des Art. 5 der \textit{Jubiläumsschrift}: »Im Grunde gehört aber der eigentliche Inhalt
der ganzen Argumentation einem höheren von Räumlichem unabhängigen
Gebiete der allgemeinen abstrakten Grössenlehre an, dessen Gegenstand die
nach der Stetigkeit zusammenhängenden Grössenkombinationen sind, einem
Gebiete, welches zur Zeit noch wenig angebauet ist, und in welchem man
sich auch nicht bewegen kann, ohne eine von räumlichen Bildern entlehnte
Sprache.«
\par
Wir werden nun im Folgenden in den Nummern 1 und 2 zunächst die
Tatsachen beweisen, die sich auf den Verlauf der Kurven \(\Re f(z)=0,\ \Im f(z)=0\)
ausserhalb des Kreises \(|z|=R\) beziehen. Es werden dabei im Wesentlichen
nur die Überlegungen der betreffenden Teile der beiden \textsc{Gauss}schen Abhandlungen
in modernerer Weise ausgeführt. In den Nummern 3 bis 6 beweisen
wir sodann die hier in Betracht kommenden Tatsachen über den Gesamtverlauf
der Kurven \(\Re f(z)=0,\ \Im f(z)=0\) auch im Innern des Kreises \(|z|=R\)
und schliessen daran den Beweis der Wurzelexistenz im Sinne der \textsc{Gauss}schen
Methode an. Dabei werden zugleich einige von \textsc{Gauss} ohne Beweis angegebene
Tatsachen über die Kurven \(\Re f(z)=0,\ \Im f(z)=0\) hergeleitet. Der Beweis des
oben zuletzt angegebenen Satzes wird an anderer Stelle\footnote[1]{Im \textsc{Crelle}schen Journal f. d. r. u. a. M. Bd. 170.} veröffentlicht werden.
\par
